% ============================================================
% ICÔNE OMARCHY (dessin partagé)
% ============================================================
% Icône vectorisée depuis /usr/share/omarchy/icon.png : le dessin est une
% grille de 15 x 15 cases de 20 px, chaque case étant uniformément pleine ou
% vide. On la redessine donc exactement, en 13 rectangles, sans dépendre d'un
% fichier externe : l'icône suit la couleur du CV et reste nette à toute taille.
%
% Les rectangles sont dans UN SEUL \fill : un unique remplissage évite les
% filets blancs que certains lecteurs PDF affichent entre deux aplats adjacents.
%
% \omarchyicon{<côté>}{<couleur>} produit le tikzpicture seul, sans placement :
% c'est à l'appelant de le positionner (dans le flux pour la sidebar, en fond de
% page pour la variante une colonne).
\newcommand{\omarchyicon}[2]{%
  \begin{tikzpicture}[x=\dimexpr#1/15\relax, y=\dimexpr#1/15\relax]
    % Une unité = une case de la grille.
    \fill[#2]
        (0,14) rectangle (15,15)
        (0,0) rectangle (1,14)
        (7,12) rectangle (8,14)
        (14,0) rectangle (15,14)
        (2,12) rectangle (7,13)
        (11,12) rectangle (13,13)
        (2,2) rectangle (3,12)
        (12,2) rectangle (13,12)
        (1,7) rectangle (2,8)
        (3,2) rectangle (12,3)
        (7,0) rectangle (8,2)
        (1,0) rectangle (7,1)
        (9,0) rectangle (14,1);
  \end{tikzpicture}%
}
